\section{Introducción: por qué importa esta entrevista}

Esta entrevista reúne en una misma conversación un conjunto de palabras clave de la industria de los grandes modelos a inicios de 2026: OpenClaw, Agent, postentrenamiento, RL scaling, MiMo-V2, orquestación de múltiples modelos, modelo base de 1T, asignación de GPU, igualdad organizacional, el entorno importa más que la experiencia. La invitada, Luo Fuli (罗福莉), trabajó en Alibaba DAMO Academy (阿里达摩院) y en DeepSeek, y actualmente dirige el equipo de grandes modelos de Xiaomi. El valor de la entrevista no está solo en «qué tan bien se desempeña cierto modelo», sino en que muestra cómo un AI Lab percibe un cambio de paradigma y reorganiza su forma de investigar.

\begin{importantbox}{Tesis central de este episodio}
El juicio de Luo Fuli es que la guerra de los grandes modelos entra en su segundo acto: de la era del Chat, dominada por el pre-train, a la era del Agent, dominada por el post-train. En adelante, la competencia no dependerá solo del modelo base, sino también del framework de Agent, la RL infra, los entornos de tareas, el costo y la velocidad, la agilidad organizacional y la asignación de capacidad de cómputo.
\end{importantbox}

\begin{knowledgebox}{Nota sobre la estrategia visual}
Este video es una entrevista con plano fijo, sin slides didácticas, pizarrón ni demostraciones de producto. Conforme al estándar de pódcast de este repositorio, el cuerpo del texto no repite fotogramas de las personas; la portada sirve para identificar la fuente, y el contenido técnico se presenta mediante diagramas conceptuales y tablas.
\end{knowledgebox}

\subsection{Resumen de la sección}

Este episodio debe leerse como material oral de un «manual de operación de un AI Lab en la era del Agent»: no trata solo de modelos, sino de cómo cambian en conjunto los modelos, los frameworks, el postentrenamiento, la capacidad de cómputo y la organización.
